%-------------------------------------------------------------
% CV DE BASE - FRANÇAIS (modèle vierge)
% Configuration (une seule fois) : l'agent remplace chaque [[MOI: ...]]
% avec le vrai CV de l'utilisateur, et ajuste le nombre d'entrées et de
% puces à son vrai parcours. Le CV doit tenir sur 1 page.
% Ensuite, ne jamais modifier ce fichier pour une offre : le skill le
% copie dans candidatures/<dossier>/cv.tex et modifie la copie.
% FORMAT FIGÉ : ne rien changer entre \documentclass et
% \begin{document} (polices, marges, espacements, commandes).
%-------------------------------------------------------------
\documentclass[letterpaper,10pt]{article}

\usepackage[letterpaper,left=0.6in,right=0.6in,top=0.42in,bottom=0.4in]{geometry}
\usepackage[T1]{fontenc}
\usepackage[utf8]{inputenc}
\usepackage[french]{babel}
\frenchsetup{StandardLists=true}
\usepackage{enumitem}
\usepackage[hidelinks]{hyperref}
\usepackage{titlesec}
\input{glyphtounicode}
\pdfgentounicode=1

\pagestyle{empty}
\setlength{\parindent}{0pt}

\titleformat{\section}{\large\scshape}{}{0em}{}[\titlerule]
\titlespacing*{\section}{0pt}{15pt}{6pt}

\setlist[itemize]{leftmargin=1.15em, itemsep=2pt, topsep=5pt, parsep=0pt, partopsep=0pt}

%----------COMMANDES----------
% \experience{Entreprise}{Dates}{Poste}{Lieu}
\newcommand{\experience}[4]{\textbf{#1}\hfill #2\\ \textit{#3}\hfill\textit{#4}}
% \projet{Nom -- Description}{(technos ou contexte)}{Dates}
\newcommand{\projet}[3]{\textbf{#1} {\small #2}\hfill #3}
% espace entre deux entrées
\newcommand{\entre}{\vspace{8pt}}

%=============================================================
\begin{document}

%----------EN-TÊTE----------
\begin{center}
    {\Huge \textbf{[[MOI: Prénom Nom]]}} \\[3pt]
    \href{mailto:[[MOI: courriel]]}{[[MOI: courriel]]}\enspace|\enspace{}[[MOI: téléphone]]\enspace|\enspace\href{[[MOI: URL LinkedIn]]}{LinkedIn}\enspace|\enspace\href{[[MOI: URL GitHub]]}{GitHub}
\end{center}
\vspace{-6pt}

%----------EXPÉRIENCE----------
\section{Expérience}

\experience{[[MOI: Entreprise 1]]}{[[MOI: Mois AAAA -- Mois AAAA]]}{[[MOI: Poste]]}{[[MOI: Ville, Province]]}
\begin{itemize}
  \item [[MOI: Verbe d'action + ce que tu as fait + résultat (chiffré si possible)]]
  \item [[MOI: Réalisation 2]]
  \item [[MOI: Réalisation 3]]
  \item [[MOI: Réalisation 4]]
\end{itemize}
\entre

\experience{[[MOI: Entreprise ou association 2]]}{[[MOI: Mois AAAA -- présent]]}{[[MOI: Poste]]}{[[MOI: Ville, Province]]}
\begin{itemize}
  \item [[MOI: Réalisation 1]]
  \item [[MOI: Réalisation 2]]
  \item [[MOI: Réalisation 3]]
\end{itemize}

%----------PROJETS----------
\section{Projets}

\projet{[[MOI: Projet 1 -- Description courte]]}{([[MOI: technos principales]])}{[[MOI: Mois AAAA -- Mois AAAA]]}
\begin{itemize}
  \item [[MOI: Ce que fait le projet et ton rôle]]
  \item [[MOI: Défi technique résolu et comment]]
  \item [[MOI: Résultat ou impact]]
\end{itemize}
\entre

\projet{[[MOI: Projet 2 -- Description courte]]}{([[MOI: technos ou contexte, ex. hackathon, équipe de 3]])}{[[MOI: Mois AAAA]]}
\begin{itemize}
  \item [[MOI: Réalisation 1]]
  \item [[MOI: Réalisation 2]]
  \item [[MOI: Réalisation 3]]
\end{itemize}
\entre

\projet{[[MOI: Projet 3 -- Description courte]]}{}{[[MOI: Mois AAAA -- Mois AAAA]]}
\begin{itemize}
  \item [[MOI: Réalisation 1]]
  \item [[MOI: Réalisation 2]]
  \item [[MOI: Réalisation 3]]
\end{itemize}

%----------FORMATION----------
\section{Formation}

\textbf{[[MOI: Établissement]]}\hfill [[MOI: Ville, Province]]\\
\textit{[[MOI: Programme, ex. Baccalauréat en génie logiciel]]}\hfill\textit{[[MOI: Diplomation prévue en mois AAAA]]}\\
[[MOI: Crédits complétés, moyenne, prix… ou supprimer cette ligne]]

%----------COMPÉTENCES TECHNIQUES----------
\section{Compétences techniques}

\textbf{Langages :} [[MOI: Python, Java, ...]]\\[2.5pt]
\textbf{Frameworks :} [[MOI: React, Node.js, ...]]\\[2.5pt]
\textbf{Outils :} [[MOI: Git, Linux, Docker, ...]]\\[2.5pt]
\textbf{Méthodologies :} [[MOI: Agile/Scrum, tests unitaires, ...]]

\end{document}
